% Generated draft: requires mathematical and source review.
\documentclass[11pt]{article}
\usepackage[margin=25mm]{geometry}
\usepackage{fontspec}
\setmainfont{Latin Modern Roman}
\usepackage{amsmath,amssymb,mathtools,unicode-math}
\setmathfont{Latin Modern Math}
\usepackage{enumitem,needspace,xurl,hyperref}
\setlength{\parindent}{0pt}
\setlength{\parskip}{5pt}
\setlength{\emergencystretch}{3em}
\newcommand{\N}{\mathbb{N}}
\newcommand{\Z}{\mathbb{Z}}
\newcommand{\Q}{\mathbb{Q}}
\newcommand{\R}{\mathbb{R}}
\newcommand{\C}{\mathbb{C}}
\newcommand{\F}{\mathbb{F}}
\newcommand{\A}{\mathbb{A}}
\newcommand{\PP}{\mathbb{P}}
\newcommand{\E}{\mathbb{E}}
\newcommand{\eps}{\varepsilon}
\newcommand{\dd}{\,\mathrm{d}}
\begin{document}
\section*{UnsolvedMath ID: 1873}
\textbf{Source problem code:} catalogue entry
\subsection{Definitions and mathematical statement}
The requested object is the least positive odd integer not expressible as $\pm p^a\pm2^b$, with $p$ an odd prime. The record does not state the allowed exponent sets. To expose this missing datum, let $E_a$ and $E_b$ denote chosen sets of allowed exponents. Under the usual independent-sign reading, define
\[\mathcal R(E_a,E_b)=\{\epsilon_1p^a+\epsilon_2 2^b:p\text{ is an odd prime},\ a\in E_a,\ b\in E_b,\ \epsilon_1,\epsilon_2\in\{-1,1\}\}.\]
The target, once conventions have been fixed, is
\[\min\bigl(\{1,3,5,\ldots\}\setminus\mathcal R(E_a,E_b)\bigr),\]
provided this set is nonempty. In particular, allowing exponent zero, requiring strictly positive exponents, or tying the two signs together can change the representable set. Neither an exponent convention nor a different sign convention is silently selected here. The question requires excluding every allowed representation of a candidate and verifying that each smaller positive odd integer has an allowed representation.

The supplied background attributes to Cohen and Selfridge the observation that the requested number exceeds $2^{18}$, but does not supply its exact exponent/sign conventions. The literature triage checked 2026-08-17 reports infinitely many nonrepresentable odd integers in related forms, including an arithmetic progression that is neither a sum nor a difference of a power of two and a prime, and later covering-congruence constructions. Such related forms must be distinguished from the full displayed expression allowing $p^a$: the review explicitly says the least number for this exact record was not verified and requests a dedicated formulation and computational/status audit. The registry labels the problem open; the dated review labels it partially solved / PARTIAL-PROGRESS. No exact least value or exact-form existence result is inferred from the related constructions.
\subsection{Short English statement}
Find the first positive odd integer that cannot be written as a signed sum of an odd-prime power and a power of two. The exponent ranges and sign conventions must be fixed before identifying the first exception; the supplied results concern related nonrepresentation problems.
\subsection{Sources}
\begin{itemize}
\item Local UnsolvedMath registry, permanent ID 1873, source code catalogue entry, ``Cohen-Selfridge Problem on $\pm p^a\pm2^b$''; supplied statement and background.
\item Richard Guy, \emph{Unsolved Problems in Number Theory}, Section A19, as attributed by the source; no edition or page is supplied.
\item F. Cohen and J. L. Selfridge, ``Not Every Number is the Sum or Difference of Two Prime Powers,'' \emph{Mathematics of Computation} 29 (1975), 79--81: \url{https://www.ams.org/journals/mcom/1975-29-129/S0025-5718-1975-0376583-0/S0025-5718-1975-0376583-0.pdf}. Supplied primary reference for related arithmetic-progression nonrepresentation results.
\item Z.-W. Sun, ``On integers not of the form plus/minus $p^a$ plus/minus $q^b$,'' \emph{Journal of Number Theory} 128 (2008), 997--1002: \url{https://doi.org/10.1016/j.jnt.2007.06.003}. Supplied primary reference for stronger covering-congruence results.
\item Supplied literature triage checked 2026-08-17, PARTIAL-PROGRESS; it preserves the wording, calls for an exact-form audit, and cautions that dated triage does not exclude later or unindexed work.
\end{itemize}
\end{document}
